% UNREGISTERED D.73 candidate. Await root and independent audits.
\subsection{Class-two cusp rigidity with any number of punctures}

\begin{statement}{D.73}{Draft}{Generic class-two cusp fibres on every affine hyperbolic curve}
Let \(U=X-D\) be a smooth geometrically connected affine
hyperbolic curve over a number field \(k\). Assume only that
\(D\) contains one \(k\)-rational point \(d\), and choose a
nonzero rational tangent there. No condition is imposed on
the number of other geometric punctures or their fields of
definition. Use this actual cusp section \(s_0\) as reference,
and put
\[
 P_\ell=\pi_1(U_{\bar k})^{(\ell)},\qquad
 T_{2,\ell}=P_\ell/\gamma_3P_\ell,\qquad
 {\cal K}_{k,\ell}=H^1(G_k,\mathbb Z_\ell(1)).
                                                        \tag{D.73.1}
\]
The chosen cusp packet injects into
\(H^1(G_k,T_{2,\ell})\), at every prime \(\ell\), including
two. In the intrinsic topology on \({\cal K}_{k,\ell}\)
specified in D.70 there is a dense \(G_\delta\) of
infinitely ramified parameters whose actual full pro-\(\ell\)
section fibre over this class-two image is the singleton
chosen cusp class.

The product generic locus in
\(\prod_\ell{\cal K}_{k,\ell}\) has the same singleton
property for the actual full profinite class-two fibre,
using precisely D.24 and D.56. Every coefficient-prime
component may simultaneously be infinitely ramified.
The compactness theorem D.72 then gives finite nilpotent
failure for every other class at a fixed finite stage,
uniformly outside any neighborhood of the cusp in the
fixed-class-two fibre. The all-prime version of this
last finite-failure assertion concerns the actual
pronilpotent tower, exactly as in D.72.

This includes curves with one or two geometric cusps.
For a one-cusp curve of genus greater than one it does
not identify the entire zero-abelian class-two fibre.
For a two-cusp curve it makes no singleton abelian
fibre claim. The exact zero-abelian parametrization
of D.71 remains specific to the once-punctured elliptic
curve. No first finite obstruction is computed.
\end{statement}

\begin{proof}
Fix \(\ell\), write \(r=|D(\bar k)|\), \(P=P_\ell\),
\(T_2=T_{2,\ell}\), \(V=P^{\rm ab}\), and
\(L_2=\gamma_2P/\gamma_3P=\bigwedge^2V\).
All these integral identifications use the actual
free pro-\(\ell\) group, with no logarithmic
coordinates. Let \(c\) generate the chosen
peripheral subgroup \(I_c\), so that
\[
                   \sigma_0(g)c=c^{\chi_\ell(g)}.       \tag{D.73.2}
\]
Choose a finite set \(S\) containing the places above
\(\ell\), the archimedean places, the bad places of
the marked pair, and the places where the chosen
tangent is not integral and nonzero.

\paragraph{Injection of the chosen packet.}
At a good \(v\notin S\), the Frobenius eigenvalues
on \(V\) have complex absolute values \(N(v)^{1/2}\)
or \(N(v)\), as proved in D.72. Those on \(L_2\)
have absolute values among \(N(v),N(v)^{3/2},N(v)^2\).
Thus neither lattice has a nonzero local fixed
element, and restriction to one such place gives
the corresponding global assertion.

If \(r\ge2\), the image of \(c\) in the boundary
part of \(V\) is primitive, and \(V/I_c\) is
torsion-free. Its boundary/proper sequence is
the quotient of (D.72.6) by this one boundary
line. Its Frobenius eigenvalues still have
positive weights, so
\((V/I_c)^{G_{k_v}}=(V/I_c)^{G_k}=0\).
The integral coefficient exact sequence gives
an injection
\[
 H^1(G,I_c)\lhook\joinrel\longrightarrow H^1(G,V),
                 \quad G=G_k\text{ or }G_{k_v}.        \tag{D.73.3}
\]
Consequently two parameters in the chosen
packet cannot have equal class-two classes
unless they are equal.

If \(r=1\), hyperbolicity gives \(g=g(X)\ge1\)
and \(V=T_\ell J_X\). In an actual surface
basis \(a_1,b_1,\ldots,a_g,b_g\), the boundary
word is the product of the commutators, up
to inversion. Its degree-two class is
\[
                  \Omega=\sum_{i=1}^g[a_i,b_i]
                        \quad\text{in }\bigwedge^2V.  \tag{D.73.4}
\]
It is primitive, including at two, and its
line \(I=\mathbb Z_\ell\Omega\) is the actual
image of \(I_c\simeq\mathbb Z_\ell(1)\).
The quotient \(L_2/I\) is torsion-free and
has only positive Frobenius weight, hence
zero local and global fixed subgroup.
Therefore \(H^1(G,I)\to H^1(G,L_2)\) is
injective. The map from this latter group
to the zero-abelian fibre of \(H^1(G,T_2)\)
is injective too: a conjugator between
two central cocycles has fixed image in
\(V\), hence belongs to \(L_2\), where
its relation is the usual coboundary
relation. This proves the chosen cusp
injection globally and locally for \(r=1\).
It does not say \(I=L_2\) unless \(g=1\).

More generally, for any section action
on \(T_2\), a fixed element first has
zero image in \(V\) and then lies in
the zero fixed subgroup of \(L_2\).
Inner twisting changes neither graded
action. We have proved
\[
               T_2^{G_k,\alpha}
                 =T_2^{G_{k_v},\alpha}=1
                              \quad(v\notin S)         \tag{D.73.5}
\]
for every actual section action \(\alpha\).

\paragraph{Peripheral centralizers and class-two separation.}
Every boundary generator has a primitive
first nonzero lower-central class in degree
at most two: degree one if \(r\ge2\), and
(D.73.4) if \(r=1\). It therefore has no
\(\ell\)-th root in \(P\), by torsion-freeness
of \(V\) and of the second graded lattice.
The closed two-generated subgroup argument
of D.71, proved there from
\cite[Theorem 3.2]{DE_Lubotzky1982}, applies
to every free finite-rank pro-\(\ell\) group.
For \(b\ne0\) it gives
\[
       C_P(c^b)=N_P(c^{b\mathbb Z_\ell})
                 =\langle c\rangle,\qquad
       \gamma_3P\cap\langle c\rangle=1.                \tag{D.73.6}
\]
For the normalizer equality, conjugation
acts trivially on that first nonzero
graded class, forcing its unit exponent
on \(c^b\) to be one. This also proves
the equalities for any other cusp.

For \(r\ge3\), two distinct geometric
peripheral lines have zero intersection
in \(V\), by the boundary permutation
module modulo its diagonal relation,
as checked in D.70. Nonzero powers of
different cusps therefore cannot be
conjugate in \(T_2\). For \(r=1\) there
is no other geometric cusp.

The new case is \(r=2\), when hyperbolicity
forces \(g\ge1\). Choose actual surface
generators with boundary words \(c_1,c_2\)
and relation
\[
          \left(\prod_{i=1}^g[a_i,b_i]\right)c_1c_2=1.
                                                        \tag{D.73.7}
\]
The images of the two boundary words
in \(V\) are opposite primitive vectors.
In \(T_2\), their exact relation is
\[
                    c_2=c_1^{-1}\Omega^{-1},
       \qquad\Omega=\sum_i[a_i,b_i]\in L_2.            \tag{D.73.8}
\]
Here additive notation is used for the
central lattice when writing \(\Omega\).
Let \(\overline V=T_\ell J_X\).
The canonical linear map
\[
              \bigwedge^2V\longrightarrow
                              \bigwedge^2\overline V  \tag{D.73.9}
\]
sends every commutator \([h,c_1]\) to
zero, and sends \(\Omega\) to the
nonzero primitive symplectic bivector
of \(\overline V\). This is the
exterior-square map of lattices; it
is not the class-two proper fundamental
group quotient, which would kill the
surface relation.

Suppose \(c_1^a=h c_2^b h^{-1}\) in
\(T_2\), with \(a,b\in\mathbb Z_\ell\).
Abelianization gives \(a=-b\).
Use (D.73.8) and cancel this common
power of \(c_1\). The remaining
central equality says
\[
                        b\Omega\in[c_1,V].            \tag{D.73.10}
\]
Applying (D.73.9) yields \(b\overline\Omega=0\).
Torsion-freeness and primitivity force
\(b=0\), and then \(a=0\).
Thus distinct geometric cusp inertia
groups are separated by conjugacy
of their nonzero powers already in
class two, also in the two-cusp case.
All comparisons may be made after a
finite unramified splitting extension
at a good place; nonzero inertia and
the geometric calculation are unchanged.

\paragraph{The exact local comparison.}
Use the actual continuous Kummer
representatives on each fixed torsion
component of \({\cal K}_{k,\ell}\)
constructed in D.70. For a parameter
\(a\), put
\[
 \sigma_a(g)=\operatorname{Int}(c^{a(g)})\sigma_0(g),
                      \qquad H=\gamma_3P.              \tag{D.73.11}
\]
A full pro-\(\ell\) section with its
prescribed class-two class can be
conjugated to have exactly the chosen
class-two cocycle. Its difference
from the cusp section is an actual
\(H\)-cocycle. By (D.73.5), a
conjugator between two exact
representatives lies in \(H\).
Consequently the fibre is exactly
\[
                         H^1(G_k,H;\sigma_a).          \tag{D.73.12}
\]

At \(v\notin S\), write \(a_v\) for
the valuation coordinate of \(a\).
If \(a_v\ne0\), the prescribed
class-two inertia is nonzero.
Hoshi's actual ordinary pro-\(\ell\)
specialization result
\cite[Appendix A, A.3, Remark A.3.1 and Proposition A.6]
{DE_Hoshi2012}
says that the section's nonzero
specialized inertia is peripheral.
The class-two separation just proved
forces it to be the chosen geometric
cusp, even if another cusp only
splits after an unramified extension.
The centralizer and normalizer
calculation (D.73.6), together with
the cartesian specialization square,
puts the entire local section in the
actual chosen cusp decomposition
group, up to geometric conjugacy.
The packet injection (D.73.3), or
the one-cusp argument following
(D.73.4), identifies its Kummer
parameter with \(a|_{G_{k_v}}\).
After this cusp adjustment the
conjugator has fixed image in
\(T_2\), so (D.73.5) puts it in
\(H\). Thus every global cocycle
in (D.73.12) is an integral local
\(H\)-coboundary at such a place.

If \(a_v=0\), a nonzero specialized
inertia group would be peripheral
and would retain a nonzero
class-two image, contradicting
the prescribed image. The same
cartesian square therefore makes
the actual inertia image zero.
The reference tangent is unramified
at this good place. Hence the
actual difference cocycle satisfies
\[
 \begin{cases}
 z|_{G_{k_v}}\text{ is an }H\text{-coboundary},
                                             &a_v\ne0,\\
 z|_{I_v}=1,                                &a_v=0.
 \end{cases}                                             \tag{D.73.13}
\]
The first assertion is over \(k_v\)
itself: the geometric comparison
over a splitting extension has
identified the chosen rational
cusp, whose full local arithmetic
decomposition group is defined
over \(k_v\). No descent of an
arbitrary geometric conjugator
from a field extension is used.

\paragraph{Finite support, witnesses, and exact descent.}
We spell out why the same Baire
argument applies to \(H=\gamma_3P\).
Let \(Q_m=P/\Phi^mP\); these are
finite characteristic quotients
of exponent dividing \(\ell^m\),
and their kernels form a
neighborhood basis. Let \(H_m\)
be the image of \(H\). Define
\[
         S_m(a)=S\cup\{v:a_v\notin\ell^m\mathbb Z_\ell\}.
                                                        \tag{D.73.14}
\]
This is finite. Outside this set,
the action \(\sigma_a\) on \(Q_m\)
is unramified. Equation (D.73.13)
makes every actual \(H\)-cocycle
unramified there modulo \(H_m\):
in the nonzero case the integral
local coboundary has trivial
inertia after reduction. Therefore
all actual reductions occur in
the finite set
\[
             Z^1(G_{k,S_m(a)},H_m;\sigma_a).            \tag{D.73.15}
\]
Finiteness is the finite-image
Hermite--Minkowski argument used
in D.70; it applies to the
associated finite semidirect
representation, not just to
abelian cohomology. Both the
finite action and this finite
support are locally constant
at sufficiently deep precision
on each fixed torsion component.

Here the action used for the
witness is precisely
\(\rho_m:G_k\to\operatorname{Aut}(H_m)\);
freezing the action on \(Q_m\)
also freezes this induced action.
If a cocycle in (D.73.15) is
nontrivial on \(\ker\rho_m\), its affine
image contains an actual element
\((z(\delta),1)\) with
\(z(\delta)\ne1\). Chebotarev
\cite[V, Theorem 3.23]{PBC_MilneCFT}
supplies fresh good primes whose
Frobenius has this action-trivial
but cocycle-nontrivial image.
The perturbations of D.70,
using elements \(\alpha_v\) with
\(\operatorname{div}(\alpha_v)=h v\)
for the class number \(h\),
give a nonzero new \(a_v\) while
preserving all previous finite
data. Choose depth beyond
\(m\) and beyond the prescribed
neighborhood, and avoid the
finite support where the old
valuation has depth at most
\(M+v_\ell(h)\). Then adding
\(\ell^M\kappa(\alpha_v)\)
has exactly nonzero valuation
of depth \(M+v_\ell(h)\)
at this fresh place, without
cancellation. The fixed torsion
component, including at two,
does not change.

The first case of (D.73.13)
now makes that Frobenius value
zero modulo \(H_m\), since its
finite action is the identity.
This contradicts the chosen
finite affine witness. As in
D.70, the finite lists produce
dense open conditions at each
precision eliminating all such
exceptional cocycles. Their
countable intersection is a
dense \(G_\delta\) in the
complete Kummer parameter space;
intersect also the dense open
conditions supplying infinitely
many nonzero valuation primes.

For a parameter in this locus,
every actual \(H\)-cocycle
factors through the actual
compact image
\(\mathcal I_a\) of \(G_k\)
in \(\operatorname{Aut}(H)\),
as detected by these finite
characteristic quotients.
Choose one nonzero valuation
place. Its inertia image is
\[
 N=\{\operatorname{Int}(c^{a_v t})|_H:
                             t\in\mathbb Z_\ell\}.
                                                        \tag{D.73.16}
\]
It is normal in \(\mathcal I_a\)
by (D.73.2). By (D.73.6),
\(H^N=1\), for primitive or
nonprimitive \(a_v\).
Its restriction is a coboundary
by (D.73.13); change the
global representative to kill
this restriction. For \(n\in N\)
and \(g\in\mathcal I_a\), the
two cocycle expansions of
\(ng=g(g^{-1}ng)\) show that
the value at \(g\) is fixed
by every \(n\). It is therefore
one. This proves the singleton
in (D.73.12).

Finally take the product of
the residual loci over all
coefficient primes. The
topology and actual full cusp
existence are precisely those
proved in D.70. If an actual
full section has the given
class-two class, its full
pro-\(\ell\) image is that
fixed cusp image for every
\(\ell\). D.24/D.56 then
compares the two actual full
sections and makes them equal.
Apply D.72 to this fixed
class-two cusp to obtain the
finite-stage and compact
nested-image assertions.
There is no claim that the
entire zero-abelian fibre has
been computed outside D.71's
elliptic case, or that
primewise geometric quotients
are jointly faithful.
\end{proof}
